\section{Introducción}
El desarrollo de software contemporáneo enfrenta retos significativos, derivados de la necesidad de construir sistemas escalables, seguros y mantenibles. En el ámbito del desarrollo backend, la comunicación eficiente entre el cliente y el servidor se ha convertido en un pilar, popularizando los servicios Web RESTful por su simplicidad y eficiencia frente a tecnologías más complejas como SOAP.
Para abordar la creciente complejidad del código y asegurar su calidad a largo plazo, la comunidad técnica ha promovido la adopción de estructuras de diseño y paradigmas que fomentan la separación de responsabilidades y la reutilización. En este sentido, la elección del framework backend es vital, donde Spring Boot y Node.js se presentan como contendientes robustos. De forma complementaria, la aplicación de patrones arquitectónicos como el Modelo-Vista-Controlador (MVC) busca imponer modularidad, coherencia y claridad, facilitando las pruebas unitarias y la evolución del sistema.
En este contexto, el presente trabajo busca analizar y aplicar estas prácticas a nivel técnico, con el objetivo de demostrar que la integración de frameworks backend maduros (como Spring Boot y su afinidad con MVC) junto con la aplicación de la metodología Scrum, favorece la creación de sistemas más robustos, seguros y alineados con los requisitos funcionales del proyecto.